
\subsection{Versionnement}

\begin{frame}{Versionnement}

Toute transaction $T$ en cours a deux choix à chaque instant

\begin{redbullet}
  \item Valider les maj effectuées avec \texttt{commit}.
  \item Les annuler avec \texttt{rollback}.
\end{redbullet}

\vfill

Le SGBD maintient, pendant l'exécution de $T$, deux versions des nuplets mises à jour :

\begin{redbullet}
   \item une version du nuplet \alert{après} la mise à jour ;
    \item une version du nuplet \alert{avant} la mise à jour.
\end{redbullet}

\vfill
\begin{blocgauche}
On parle \alert{d'image après} et \alert{d'image avant} pour ces deux versions.
\end{blocgauche}
\end{frame}

\begin{frame}{Isolation basée sur les versions}

Soit deux transactions $T$ et $T'$. 

\begin{redbullet}

    \item  Chaque fois que $T$ met à jour un nuplet, la version courante est 
       copiée dans \alert{l’image avant}, puis remplacée par la valeur de \alert{l’image après} fournie par $T$.
        
     \item Quand $T$ lit des nuplets, 
        le système doit lire dans l’image après pour assurer une vision cohérente de la base
      
    \item En revanche,  une autre transaction $T'$ 
        doit lire  dans \alert{l’image avant} pour éviter les effets de lectures sales.
\end{redbullet}


\vfill
\begin{blocgauche}
\alert{Combien d'images avant / après}? Une seule mise à jour est autorisée à la fois, donc
une seule paire (image avant, image après). 
\end{blocgauche}
\end{frame}


\begin{frame}{Illustration}

Le nuplet $e_2$ est en cours de modification.


\medskip

\figSlideWithSize{read-committed}{13cm}

\vfill

$T_1$ et $T_2$ ne lisent pas la même version. 


\end{frame}


\begin{frame}{Les lectures ne sont pas répétables}

Le mode précédent est correct pour le niveau \texttt{read committed}.


\vfill


\figSlideWithSize{unrepeatable-read}{13cm}


\vfill

\begin{blocgauche}
Mais les lectures ne sont pas répétables: $T_2$ lit la dernière version validée.
\end{blocgauche}
\end{frame}


\begin{frame}{Image avant et lectures répétables}

Les images avant constituent un "cliché" de la base pris à un moment donné.

\vfill

\alert{Principe des lectures répétables}: 
 toute transaction ne lit que dans le cliché valide au moment où elle a débuté.
 
\vfill

\alert{Estampillage}: 

\begin{redbullet}
   \item Toute transaction $T_i$ est ``estampillée'' par le moment $\tau_i$
     où elle a commencé. 
     \item Chaque nuplet $e$ est marqué par le moment $\tau_e$
         de sa validation (\texttt{commit}).
      \item $T_i$ ne peut lire que les images avant $e$  dont l'estampille $\tau_e < \tau_i$.   
\end{redbullet}

\begin{blocgauche}
Ce mécanisme de \alert{versionnement} assure les lectures cohérentes.
\end{blocgauche}

\end{frame}


\begin{frame}{Lectures répétables (niveau \texttt{repeatable read}) }

Une transaction ne lit que
les nuplets dont l'estampille est inférieure à la sienne. 

\vfill


\figSlideWithSize{repeatable-read}{13cm}


\vfill
\begin{blocgauche}\alert{Combien de versions?} Il faut maintenir le ``cliché'' de la base pris
  pour la transaction la plus ancienne en cours d'exécution.
\end{blocgauche}

\end{frame}


%%%%%%%%%%%%%%%%%%%%%%%%%%%%%%%%%%%%%%%%%%%%%%%%%%%%%%%
\begin{frame}{À retenir}

Notion \alert{d'image avant} et \alert{image après}
\begin{redbullet}
  \item \alert{Image après}\,: les dernières versions de chaque nuplet; certaines
      ne sont pas encore validées.
       
  \item \alert{Image avant}: toutes les versions antérieures
\end{redbullet}

\vfill

Le multi-versions assure les niveaux \texttt{read committed} et \texttt{repeatable read}.

\begin{redbullet}
  \item \texttt{read committed} : il suffit de conserver une version dans l'image avant
  \item \texttt{repeatable read} : il faut conserver les versions immédiatement antérieure
    au début de la transaction (``cliché'' = \textit{snapshot})
\end{redbullet}

\vfill

Niveau plus élevé = performances affectées
\end{frame}
